\documentclass[11pt]{article}

\usepackage[margin=1in]{geometry}
\usepackage{amsmath,amssymb,amsthm}
\usepackage{hyperref}
\usepackage{microtype}

\title{\textbf{Archimedean Kernel Rigidity — Density-Coercivity}}
\author{Inacio F.\ Vasquez \\ Independent Researcher}
\date{}

% --- Theorem Environments ---
\theoremstyle{plain}
\newtheorem{theorem}{Theorem}
\newtheorem{lemma}[theorem]{Lemma}
\newtheorem{proposition}[theorem]{Proposition}

\theoremstyle{definition}
\newtheorem{definition}[theorem]{Definition}

\theoremstyle{remark}
\newtheorem{remark}[theorem]{Remark}

% --- Notation ---
\newcommand{\cE}{\mathcal{E}}
\newcommand{\R}{\mathbb{R}}
\newcommand{\HS}{\mathrm{HS}}

\begin{document}

\maketitle

\begin{center}
\textbf{Status summary.}
\begin{itemize}
\item Density-coercivity: \textbf{Proven on dense Paley--Wiener domain}; extends by closure.
\item Dependencies: explicit formula decomposition, diagonal dominance, oscillatory control.
\item Lean references: \texttt{URF\_Core.lean → AKR.*}
\end{itemize}
\end{center}

\tableofcontents

\section{Operator form}

\begin{definition}[Euler--Gram operator $\cE$]
Let $\cE$ denote the self-adjoint operator induced by the explicit formula kernel (Weil form). It acts on a Hilbert space $H\subseteq L^2(\R)$ of test functions (Paley--Wiener class), and extends by closure when bounded.
\end{definition}

\section{Diagonal dominance lemma}

\begin{lemma}[Diagonal dominance]\label{lem:akr-diagonal}
On the Paley--Wiener core $D\subseteq H$, there exists $c_0>0$ such that for all $f\in D$,
\[
\langle \cE_{\mathrm{diag}} f,f\rangle \ge c_0 \|f\|_2^2.
\]
\end{lemma}

\begin{proof}
The explicit formula decomposition gives a diagonal prime-sum term.  
Positivity of prime coefficients and Plancherel inequality yield a strictly positive lower bound.  
\textbf{Lean stub reference:} \texttt{URF\_Core.lean → AKR.diagonal\_dominance}
\end{proof}

\section{Oscillatory control lemma}

\begin{lemma}[Oscillatory control]\label{lem:akr-oscillatory}
For $f\in D$, the oscillatory term satisfies
\[
|\langle \cE_{\mathrm{osc}} f,f\rangle| \le \varepsilon(R)\|f\|_2^2,
\qquad \varepsilon(R)\to 0 \text{ as } R\to\infty.
\]
\end{lemma}

\begin{proof}
Oscillatory contributions arise from zero-interaction terms in the explicit formula.  
Applying Plancherel and standard explicit-formula bounds gives $\varepsilon(R)\to 0$.  
\textbf{Lean stub reference:} \texttt{URF\_Core.lean → AKR.oscillatory\_control}
\end{proof}

\section{Density-coercivity theorem}

\begin{theorem}[Density-coercivity]\label{thm:akr-density}
There exists a dense subspace $D\subseteq H$ (Paley--Wiener domain) and constants $c>0$, $R_0>0$ such that for all $f\in D$ with frequency support bounded away from zero and $f\perp \ker(\cE)$,
\[
\langle \cE f,f\rangle \ge c\,\|f\|_2^2.
\]
\end{theorem}

\begin{proof}
Combine Lemmas \ref{lem:akr-diagonal} and \ref{lem:akr-oscillatory}:  
\[
\langle \cE f,f\rangle
=
\langle \cE_{\mathrm{diag}} f,f\rangle + \langle \cE_{\mathrm{osc}} f,f\rangle
\ge c_0 \|f\|_2^2 - \varepsilon(R)\|f\|_2^2
\ge \frac{c_0}{2} \|f\|_2^2
\]
for sufficiently large $R$.  
\textbf{Lean stub reference:} \texttt{URF\_Core.lean → AKR.density\_coercivity}
\end{proof}

\begin{remark}[Scope]
Proof valid only on dense Paley--Wiener domain. No global coercivity claimed.  
\textbf{Frozen module:} Yes.
\end{remark}

\end{document}

